\documentclass[12pt]{article}

% Idioma e codificação
\usepackage[brazil]{babel}
\usepackage[utf8]{inputenc}
\usepackage[T1]{fontenc}

% Matemática e listas
\usepackage{amsmath, amssymb, amsfonts}
\usepackage{enumitem}
\usepackage{tabularx}
\usepackage{booktabs}

% Layout e cores
\usepackage{geometry}
\usepackage{graphicx}
\usepackage{xcolor}

% Cores institucionais
\definecolor{maincolor}{RGB}{0,100,50}
\definecolor{forestgreen}{RGB}{0,100,50}

% Cabeçalho e rodapé
\usepackage{fancyhdr}

% Blocos coloridos
\usepackage[most]{tcolorbox}
\newtcolorbox{block}[1]{%
  colback=green!5,
  colframe=maincolor,
  coltitle=white,
  title=#1,
  fonttitle=\bfseries,
  boxrule=1.5pt,
  arc=3mm,
  left=6pt, right=6pt, top=6pt, bottom=6pt
}

% Margens
\geometry{margin=2.5cm}

% Fancyhdr
\pagestyle{fancy}
\fancyhf{}

\renewcommand{\headrulewidth}{2pt}
\renewcommand{\footrulewidth}{2pt}

\renewcommand{\headrule}{\hbox to\headwidth{%
  \color{forestgreen}\leaders\hrule height \headrulewidth\hfill}}

\renewcommand{\footrule}{\hbox to\headwidth{%
  \color{forestgreen}\leaders\hrule height \footrulewidth\hfill}}

\setlength{\headheight}{52pt}
\addtolength{\topmargin}{-20pt}

\fancyhead[L]{%
  \includegraphics[height=1.2cm]{Universidade_de_Rio_Verde.jpg}
}
\fancyhead[R]{%
  \raggedleft
  \textcolor{forestgreen}{%
    \small
    \textbf{Lista de Exercícios}\\
    Me.~Sergio Souza Novak\\
    \textit{Banco de Dados}
  }
}
\fancyfoot[C]{\textcolor{forestgreen}{\thepage}}

\begin{document}

\vspace*{0.2cm}
\begin{center}
  {\Large \textbf{Lista de Exercícios: Normalização}}\\[0.2cm]
  {\large Banco de Dados}\\[0.2cm]
  Tema: 4ª Forma Normal (4FN) --- Dependências Multivaloradas
\end{center}

\vspace{0.3cm}

\begin{block}{Instruções}
Para cada questão: (1) identifique as Dependências Multivaloradas (DMV) usando a notação $X \twoheadrightarrow Y$; (2) verifique se o determinante é superchave; (3) conclua com \textbf{Viola} ou \textbf{Não viola} a 4FN e justifique. Nas questões de projeto, apresente as novas tabelas com dados e indique as chaves (PK/FK).
\end{block}

\vspace{0.4cm}

\begin{enumerate}[label=\textbf{Questão \arabic*:}, leftmargin=*]

% ================================================================
%  BLOCO 1 — IDENTIFICAR DMV E TESTAR 4FN (Q1–Q3)
% ================================================================

  \item \textbf{(Teste 4FN -- Pesquisadores)} A tabela \texttt{PESQUISADORES} abaixo registra, para cada pesquisador, os projetos nos quais atua e os idiomas que domina. Um pesquisador pode atuar em vários projetos \textbf{e} dominar vários idiomas, de forma \textbf{independente}. Chave Primária: \textbf{\{Pesquisador, Projeto, Idioma\}}.

  \vspace{0.2cm}
  \centering
  \begin{tabular}{|l|l|l|}
    \hline
    \textbf{Pesquisador} & \textbf{Projeto} & \textbf{Idioma} \\ \hline
    Ana Souza  & IA Climática   & Inglês    \\ \hline
    Ana Souza  & IA Climática   & Espanhol  \\ \hline
    Ana Souza  & Robótica       & Inglês    \\ \hline
    Ana Souza  & Robótica       & Espanhol  \\ \hline
    Bruno Lima & Robótica       & Inglês    \\ \hline
    Bruno Lima & Robótica       & Francês   \\ \hline
    Bruno Lima & Bioinformática & Inglês    \\ \hline
    Bruno Lima & Bioinformática & Francês   \\ \hline
    Carla Dias & IA Climática   & Português \\ \hline
    Carla Dias & Bioinformática & Português \\ \hline
  \end{tabular}
  \flushleft
  \vspace{0.2cm}
  \begin{enumerate}[label=\alph*)]
    \item Identifique as Dependências Multivaloradas (DMV) existentes, usando a notação $X \twoheadrightarrow Y$.
    \item Por que a combinação entre \textit{Projeto} e \textit{Idioma} constitui um \textbf{produto cartesiano}? Calcule quantas linhas seriam necessárias se Ana tivesse 4 projetos e 3 idiomas.
    \item A tabela está na \textbf{4FN}? O determinante das DMVs é superchave? Justifique sua conclusão.
  \end{enumerate}

  \vspace{0.4cm}

  \item \textbf{(Teste 4FN -- Desenvolvedores)} A tabela \texttt{DESENVOLVEDORES} abaixo registra, para cada desenvolvedor, os sistemas em que trabalha e as linguagens de programação que domina. As duas informações são \textbf{independentes} entre si. Chave Primária: \textbf{\{Desenvolvedor, Sistema, Linguagem\}}.

  \vspace{0.2cm}
  \centering
  \begin{tabular}{|l|l|l|}
    \hline
    \textbf{Desenvolvedor} & \textbf{Sistema} & \textbf{Linguagem} \\ \hline
    Felipe Melo  & ERP Financeiro  & Java       \\ \hline
    Felipe Melo  & ERP Financeiro  & Python     \\ \hline
    Felipe Melo  & Portal RH       & Java       \\ \hline
    Felipe Melo  & Portal RH       & Python     \\ \hline
    Gisele Rocha & Portal RH       & Kotlin     \\ \hline
    Gisele Rocha & Portal RH       & JavaScript \\ \hline
    Gisele Rocha & App Mobile      & Kotlin     \\ \hline
    Gisele Rocha & App Mobile      & JavaScript \\ \hline
    Hugo Neves   & ERP Financeiro  & C\#        \\ \hline
    Hugo Neves   & App Mobile      & C\#        \\ \hline
  \end{tabular}
  \flushleft
  \vspace{0.2cm}
  \begin{enumerate}[label=\alph*)]
    \item Identifique as DMVs presentes na tabela usando a notação $X \twoheadrightarrow Y$.
    \item Explique formalmente por que as dependências $\text{Desenvolvedor} \twoheadrightarrow \text{Sistema}$ e $\text{Desenvolvedor} \twoheadrightarrow \text{Linguagem}$ são \textbf{independentes}.
    \item A tabela viola a \textbf{4FN}? Identifique a anomalia de inserção que ocorreria ao adicionar a linguagem \textit{TypeScript} para Felipe Melo.
  \end{enumerate}

  \vspace{0.4cm}

  \item \textbf{(Teste 4FN -- Professores)} A tabela \texttt{PROFESSORES\_AULAS} abaixo registra os professores, as disciplinas que lecionam e os dias da semana em que estão disponíveis. A disponibilidade de um professor é \textbf{independente} da disciplina que ele leciona. Chave Primária: \textbf{\{Professor, Disciplina, Dia\}}.

  \vspace{0.2cm}
  \centering
  \begin{tabular}{|l|l|l|}
    \hline
    \textbf{Professor} & \textbf{Disciplina} & \textbf{Dia} \\ \hline
    Dr. Almeida  & Banco de Dados & Segunda  \\ \hline
    Dr. Almeida  & Banco de Dados & Quarta   \\ \hline
    Dr. Almeida  & Redes          & Segunda  \\ \hline
    Dr. Almeida  & Redes          & Quarta   \\ \hline
    Dra. Figueiredo & Algoritmos  & Terça    \\ \hline
    Dra. Figueiredo & Algoritmos  & Quinta   \\ \hline
    Dra. Figueiredo & Estruturas  & Terça    \\ \hline
    Dra. Figueiredo & Estruturas  & Quinta   \\ \hline
    Dr. Kastner  & SO             & Segunda  \\ \hline
    Dr. Kastner  & Compiladores   & Segunda  \\ \hline
  \end{tabular}
  \flushleft
  \vspace{0.2cm}
  \begin{enumerate}[label=\alph*)]
    \item Identifique as DMVs e classifique-as como triviais ou não triviais.
    \item Aplique o \textbf{teste formal da 4FN}: para cada DMV $X \twoheadrightarrow Y$, verifique se $X$ é superchave da relação.
    \item A tabela viola a \textbf{4FN}? Que redundância concreta ocorreria se Dr. Almeida passasse a estar disponível também às Sextas?
  \end{enumerate}

% ================================================================
%  BLOCO 2 — ANÁLISE DO PRODUTO CARTESIANO (Q4)
% ================================================================

  \vspace{0.4cm}

  \item \textbf{(Análise Formal -- Produto Cartesiano)} Considere a relação \texttt{FUNCIONARIOS} que registra, para cada funcionário, os \textbf{cargos} que exerceu ao longo da carreira e os \textbf{benefícios} que recebe da empresa. Sabe-se que os conjuntos de cargos e benefícios são \textbf{independentes} entre si.

  Os conjuntos de um funcionário são:
  \[
    C(\text{Marcos}) = \{\text{Analista}, \text{Gerente}, \text{Coordenador}\}
    \qquad
    B(\text{Marcos}) = \{\text{Plano de Saúde}, \text{Vale-Alimentação}\}
  \]

  \vspace{0.2cm}
  \centering
  \begin{tabular}{|l|l|l|}
    \hline
    \textbf{Funcionário} & \textbf{Cargo} & \textbf{Benefício} \\ \hline
    Marcos & Analista     & Plano de Saúde    \\ \hline
    Marcos & Analista     & Vale-Alimentação  \\ \hline
    Marcos & Gerente      & Plano de Saúde    \\ \hline
    Marcos & Gerente      & Vale-Alimentação  \\ \hline
    Marcos & Coordenador  & Plano de Saúde    \\ \hline
    Marcos & Coordenador  & Vale-Alimentação  \\ \hline
  \end{tabular}
  \flushleft
  \vspace{0.2cm}
  \begin{enumerate}[label=\alph*)]
    \item Escreva formalmente as DMVs presentes, usando a notação $X \twoheadrightarrow Y$.
    \item Calcule e escreva explicitamente o conjunto $C(\text{Marcos}) \times B(\text{Marcos})$, mostrando todos os pares ordenados.
    \item A tabela viola a \textbf{4FN}? Demonstre que o determinante \textit{Funcionário} não é superchave na relação original.
    \item Que anomalia de \textbf{exclusão} ocorreria se o benefício ``Plano de Saúde'' fosse removido para Marcos?
  \end{enumerate}

% ================================================================
%  BLOCO 3 — PROJETO: DECOMPOR PARA A 4FN (Q5)
% ================================================================

  \vspace{0.4cm}

  \item \textbf{(Projeto Lógico -- Adequar à 4FN)} A tabela \texttt{PESSOAS} abaixo registra pessoas com seus telefones de contato e suas habilidades profissionais. As habilidades e os telefones de cada pessoa são \textbf{independentes} entre si. Chave Primária: \textbf{\{CPF, Telefone, Habilidade\}}.

  As dependências identificadas são:
  \[
    CPF \to \{\text{Nome}, \text{Cidade}, \text{Email}\}
    \qquad
    CPF \twoheadrightarrow \text{Telefone}
    \qquad
    CPF \twoheadrightarrow \text{Habilidade}
  \]

  \vspace{0.2cm}
  \centering
  \begin{tabular}{|c|l|l|c|l|l|}
    \hline
    \textbf{CPF} & \textbf{Nome} & \textbf{Cidade} & \textbf{Telefone} & \textbf{Email} & \textbf{Habilidade} \\ \hline
    001 & Ana Lima  & Goiânia     & 9001 & ana@mail.com  & Python \\ \hline
    001 & Ana Lima  & Goiânia     & 9001 & ana@mail.com  & SQL    \\ \hline
    001 & Ana Lima  & Goiânia     & 9002 & ana@mail.com  & Python \\ \hline
    001 & Ana Lima  & Goiânia     & 9002 & ana@mail.com  & SQL    \\ \hline
    002 & Bruno Costa & Rio Verde & 8001 & bruno@mail.com & Java   \\ \hline
    002 & Bruno Costa & Rio Verde & 8001 & bruno@mail.com & Kotlin \\ \hline
    002 & Bruno Costa & Rio Verde & 8002 & bruno@mail.com & Java   \\ \hline
    002 & Bruno Costa & Rio Verde & 8002 & bruno@mail.com & Kotlin \\ \hline
  \end{tabular}
  \flushleft
  \vspace{0.2cm}

  \begin{enumerate}[label=\alph*)]
    \item Comprove que a tabela está na \textbf{BCNF} (ou seja, todas as DFs $X \to Y$ têm $X$ como superchave).
    \item Identifique as duas DMVs que violam a \textbf{4FN} e explique por que o determinante $CPF$ não é superchave.
    \item Aplique o algoritmo de decomposição para a 4FN:
    \begin{enumerate}[label=\roman*.]
      \item Encontre a DMV violadora $X \twoheadrightarrow Y$ com $X$ não sendo superchave.
      \item Crie $S_1 = X \cup Y$ e $S_2 = \text{Relação} - Y$.
      \item Repita o processo até que todas as relações estejam na 4FN.
    \end{enumerate}
    \item \textbf{Reconstrua o modelo lógico} apresentando as \textbf{três} novas tabelas com:
    \begin{itemize}
      \item Nome, atributos, PK e FK de cada tabela.
      \item Os dados preenchidos conforme a tabela original (sem redundâncias).
    \end{itemize}
    \item Com a decomposição, o número de linhas armazenadas passou de \textbf{8} para quantas? Qual é o ganho em relação a uma pessoa com 5 telefones e 10 habilidades?
  \end{enumerate}

\end{enumerate}

\vspace{1cm}
\begin{center}
  \textit{``Uma relação está na 4FN se, para toda dependência multivalorada não trivial $X \twoheadrightarrow Y$, $X$ é uma superchave --- pois cada fato multivalorado deve ser armazenado em sua própria relação.''}
\end{center}

\end{document}
